\begin{frame}{Waveform ML analysis pipeline}

\small
\vspace{-0.5em}

% =========================================================
% Main pipeline
% =========================================================

\centering

\begin{tikzpicture}[
    node distance=0.38cm,
    pipelinebox/.style={
        draw=UChicagoMaroon,
        fill=PaleMaroon,
        line width=1.0pt,
        rounded corners,
        align=center,
        text width=2.0cm,
        minimum height=1.4cm,
        inner sep=4pt,
        font=\scriptsize,
        text=DarkText
    },
    flowarrow/.style={
        -{Latex[length=2mm]},
        draw=UChicagoMaroon,
        line width=1.1pt
    }
]

\node[pipelinebox] (raw) {
    \textbf{Raw waveforms}\\
    Energy OR
    Timing
    Channels
};

\node[pipelinebox, right=of raw] (prep) {
    \textbf{Preprocessing}\\
    selection + LED
};

\node[pipelinebox, right=of prep] (dataset) {
    \textbf{ML dataset}\\
    windows + targets
};

\node[pipelinebox, right=of dataset] (train) {
    \textbf{Model \\development}\\
    train + validation
};

\node[pipelinebox, right=of train] (eval) {
    \textbf{Blind evaluation}\\
    CTR
};

\draw[flowarrow] (raw) -- (prep);
\draw[flowarrow] (prep) -- (dataset);
\draw[flowarrow] (dataset) -- (train);
\draw[flowarrow] (train) -- (eval);

\end{tikzpicture}

\vspace{0.2em}

\begin{block}{Input data}
    Previous studies mainly considered energy-channel waveforms only. I applied the pipeline also on timing channel data, using the energy channel for the photopeak selection.
\end{block}
% =========================================================
% Dataset split
% =========================================================
\begin{block}{Holdout validation}

\centering

\begin{tikzpicture}[
    every node/.style={
        font=\scriptsize,
        text=DarkText
    }
]

% ---------------------------------------------------------
% Original dataset
% ---------------------------------------------------------
\draw[
    draw=UChicagoMaroon,
    fill=white,
    line width=1pt,
    rounded corners
] (0,1.55) rectangle (3.0,2.45);

\node at (1.5,2.0) {\textbf{Original dataset}};

% Arrow to split
\draw[
    -{Latex[length=2mm]},
    draw=UChicagoMaroon,
    line width=1.1pt
] (3.15,2.0) -- (4.1,2.0);

% ---------------------------------------------------------
% Development set
% ---------------------------------------------------------
\draw[
    draw=UChicagoMaroon,
    fill=PaleMaroon,
    line width=1pt,
    rounded corners
] (4.3,1.55) rectangle (7.8,2.45);

\node at (6.05,2.0) {\textbf{Development set}};

% Blind test
\draw[
    draw=UChicagoMaroon,
    fill=white,
    line width=1pt,
    rounded corners
] (8.2,1.55) rectangle (10.8,2.45);

\node at (9.5,2.0) {\textbf{Blind test}};

% Plus sign
\node[font=\normalsize\bfseries] at (8.0,2.0) {$+$};

% ---------------------------------------------------------
% Development -> Train + Validation
% ---------------------------------------------------------
\draw[
    -{Latex[length=2mm]},
    draw=UChicagoMaroon,
    line width=1.1pt
] (6.05,1.45) -- (6.05,0.95);

\draw[
    draw=UChicagoMaroon,
    fill=PaleMaroon,
    line width=1pt,
    rounded corners
] (4.4,0.0) rectangle (7.7,0.8);

\draw[
    draw=UChicagoMaroon,
    line width=0.9pt
] (6.1,0.0) -- (6.1,0.8);

\node at (5.25,0.4) {\textbf{Train}};
\node at (6.9,0.4) {\textbf{Validation}};

\end{tikzpicture}

\vspace{0.3em}

\scriptsize
The \textbf{original dataset} is split into a \textbf{development set} and a \textbf{blind test}.  
Only the \textbf{development set} is used for training and model selection.

\end{block}
\end{frame}